\section{Kaggle implementation: an unsupervised workflow without labels}
\label{sec:kaggle-mall-unsupervised}

\begin{kaggleproject}{Mall Customers --- scaling, repeated clustering, and internal evidence}
\kagglesource{https://www.kaggle.com/datasets/vjchoudhary7/customer-segmentation-tutorial-in-python}{Mall Customer Segmentation Data}
\textbf{Question.} Without a target column, how can we compare candidate values of $k$ using compactness and silhouette structure?
\end{kaggleproject}

The script scales age, annual income, and spending score, then reports K-means inertia and silhouette score for several candidate cluster counts. Neither quantity is a ground-truth accuracy measure; both are internal evidence about the representation and partition.

\textbf{Before running.} Predict the qualitative pattern you expect from the experiment before you look at the output or plot.

\begin{labmode}
Stage 4A is about evidence without labels. Decide what would count as useful evidence before looking at the reported inertia and silhouette values.
\end{labmode}

\lstinputlisting[style=mlpython]{all_experiments/lab14-mall_unsupervised_workflow.py}


\begin{labdeliverable}
Submit the scaled feature definition, inertia and silhouette evidence across candidate $k$, your chosen candidate partition, one visualization, and a decision memo stating what additional domain or stability evidence is required before treating the groups as meaningful customer segments.
\end{labdeliverable}
